\section{Setup: the prime sieve as a Markov chain on $\mathbb{Z}/m\mathbb{Z}$}
\label{sec:setup}

The whole paper rests on a careful distinction between two stationary
measures attached to the prime-sieve dynamical system. Both live on the
same finite state space $\mathbb{Z}/m\mathbb{Z}$, both are fixed by the
same transfer matrix $T_m$, and their \emph{marginals} on a single
modulus $p \mid m$ agree to numerical precision. They differ, however,
in their joint structure across distinct moduli, and that difference is
exactly the object that the two theorems of
Section~\ref{sec:proofs} address.

We assemble the relevant objects here. Subsection~\ref{ssec:gap_field}
fixes the path space $\Omega$ on which everything else is defined.
Subsection~\ref{ssec:transfer_matrix} recalls the transfer matrix~$T_m$
and the convention used for its state space (the \emph{active}
subspace, which excludes the class~$0 \bmod p$ forced to zero
by Theorem~T1).
Subsection~\ref{ssec:two_measures} introduces the two stationary
measures $\pi_m^{\mathrm{Ruelle}}$ and $\pi_m^{\mathrm{emp}}$ and
explains where the discrepancy lives.
Subsection~\ref{ssec:emergent_geometry} attaches the emergent
Bianchi-I geometry $G(\mu)$ to the spectrum of $T_m$, in the form used
throughout the rest of the paper.
Subsection~\ref{ssec:summary_objects} collects the resulting glossary.

\subsection{The gap field $\{g_n\}$ and modular residues}
\label{ssec:gap_field}

Let $p_1 = 2 < p_2 = 3 < p_3 = 5 < \cdots$ be the increasing sequence
of rational primes, and let
\begin{equation}
  g_n \;:=\; p_{n+1} - p_n \qquad (n \geq 1)
  \label{eq:gap_def}
\end{equation}
be the $n$-th \emph{prime gap}. The sequence $\{g_n\}_{n \geq 1}$ is the
unique dynamical field of Persistence
Theory~\cite{PT_MATHEMATICS_M1,PT_Monograph}: every observable of the
theory is expressible as a measurable functional of this single
sequence.

We work on the path space $\Omega = \mathbb{N}^{\mathbb{N}}$ equipped
with the cylindrical $\sigma$-algebra $\mathcal{F}$ and the canonical
shift $T : \Omega \to \Omega$, $T(\omega)_n = \omega_{n+1}$.
The map $g : \mathbb{N} \to \mathbb{N}$, $n \mapsto g_n$, picks out a
distinguished point $\omega^\ast \in \Omega$; we identify $\Omega$ with
the orbit closure of $\omega^\ast$ under $T$. Any probability measure
$\mu$ on $(\Omega, \mathcal{F})$ that is $T$-invariant will be called a
\emph{stationary measure for the gap field}.

For each modulus $m \geq 2$, the modular reduction
\begin{equation}
  r_n^{(m)} \;:=\; g_n \bmod m \;\in\; \mathbb{Z}/m\mathbb{Z}
  \label{eq:residue_def}
\end{equation}
defines a coordinate process. By the Chinese Remainder Theorem (CRT),
when $m = p_1 \cdots p_k$ is a squarefree product of distinct primes,
the joint coordinate
$\mathbf{r}_n^{(m)} := (r_n^{(p_1)}, \ldots, r_n^{(p_k)})$ identifies
$\mathbb{Z}/m\mathbb{Z}$ with $\prod_i \mathbb{Z}/p_i\mathbb{Z}$ as a
ring; we will see in Section~\ref{sec:crt_factor} that this
identification lifts to the level of transfer matrices.

For each prime $p$, write
\begin{equation}
  \mathcal{A}_p \;:=\; \sigma\!\bigl(\{r_n^{(p)} : n \in \mathbb{N}\}\bigr)
  \;\subset\; \mathcal{F}
  \label{eq:Ap_def}
\end{equation}
for the $\sigma$-algebra on $\Omega$ generated by the entire trajectory
of mod-$p$ residues. An \emph{observable on channel $p$} is, by
definition, a bounded $\mathcal{A}_p$-measurable function
$\Phi_p : \Omega \to \mathbb{R}$. The two coordinate processes
$(r_n^{(p)})_n$ and $(r_n^{(q)})_n$ for distinct $p, q$ generate the
sub-$\sigma$-algebras $\mathcal{A}_p$ and $\mathcal{A}_q$, and the
question of whether these algebras are informationally independent
under a given stationary measure is the central question of this
paper.

\subsection{The transfer matrix $T_m$}
\label{ssec:transfer_matrix}

Fix an odd prime $p \geq 3$. Theorem~T1 of the
monograph~\cite[ch.~1]{PT_Monograph} states that for $6$-rough integers
(those coprime to both $2$ and $3$), the mod-$3$ self-transitions
between successive primes are exactly forbidden:
\begin{equation}
  \Pr[r_{n+1}^{(3)} = r \mid r_n^{(3)} = r] \;=\; 0
  \qquad \text{for } r \in \{1, 2\}.
  \label{eq:T1}
\end{equation}
The mod-$p$ analogue is the assertion that residues $0 \bmod p$ are
sieved out among primes $p' > p$, so that any active state space must
exclude the class $0$. Throughout the paper we therefore work on the
\emph{active subspace}
$\mathcal{S}_p := \{1, 2, \ldots, p-1\} \subset \mathbb{Z}/p\mathbb{Z}$,
and we denote by $T_p$ the corresponding restricted transfer matrix.

\paragraph{Explicit form.}
For $p = 3$, restriction yields the canonical mod-3 matrix
\begin{equation}
  T_3 \;=\;
  \begin{pmatrix} 0 & 1 \\ 1 & 0 \end{pmatrix}
  \;=\; \mathrm{antidiag}(1, 1),
  \qquad
  \mathrm{Spec}(T_3) = \{+1, -1\}.
  \label{eq:T3_explicit}
\end{equation}
This is the unique $2 \times 2$ doubly stochastic matrix of zero
trace~\cite[Thm.~T3]{PT_Monograph}; it carries the symmetry parameter
$s = \tfrac12$ as its second eigenvalue squared, and is the engine of
the entire derivation.

For $p \geq 5$, T1 forces zero diagonal entries on $\mathcal{S}_p$, and
the remaining stochasticity (each row sums to one) plus the involutive
symmetry of the active subspace under $r \mapsto p-r$ fix
\begin{equation}
  (T_p)_{ij}
  \;=\;
  \begin{cases}
    0 & \text{if } i = j,\\[2pt]
    \displaystyle \frac{1}{p-2} & \text{if } i \neq j,
  \end{cases}
  \qquad i, j \in \mathcal{S}_p.
  \label{eq:Tp_explicit}
\end{equation}
The spectrum of $T_p$ is then
\begin{equation}
  \mathrm{Spec}(T_p) \;=\;
  \Bigl\{\, 1, \;
            \underbrace{-\tfrac{1}{p-2}, \ldots, -\tfrac{1}{p-2}}_{(p-2)\text{ times}}
       \,\Bigr\}.
  \label{eq:Tp_spectrum}
\end{equation}
The Perron eigenvalue $\lambda_0 = 1$ is simple and dominant
(spectral gap $1 - \tfrac{1}{p-2} \geq \tfrac{2}{3}$ for $p \geq 5$), and the
matrix is primitive in the sense of
Perron--Frobenius~\cite{Seneta2006}; the supporting algebraic facts
are collected in Appendix~\ref{app:perron_frobenius}.

\paragraph{Tensor structure on composite moduli.}
For $m = p_1 \cdots p_k$ a squarefree product of distinct active primes,
the CRT identification $\mathbb{Z}/m\mathbb{Z} \cong \prod_i
\mathbb{Z}/p_i\mathbb{Z}$ lifts to the tensor decomposition
\begin{equation}
  T_m \;=\; T_{p_1} \otimes T_{p_2} \otimes \cdots \otimes T_{p_k},
  \label{eq:Tm_tensor_preview}
\end{equation}
acting on the active subspace
$\mathcal{S}_m := \prod_i \mathcal{S}_{p_i}$. We defer the proof to
Theorem~\ref{thm:crt_tensor}; for now we only record that the spectrum
of $T_m$ is the multiplicative orbit
$\{\prod_i \lambda^{(p_i)}_{j_i}\}$ of the per-prime spectra, so that
the Perron eigenvalue is unique (equal to~$1$) and simple.

\paragraph{Convention warning.}
Some passages of the monograph (notably the
collider chapter, ch.~20, R58b) report
$\pi_3 = (\tfrac14, \tfrac38, \tfrac38)$ rather than the uniform
distribution on the active subspace. The two values are not in
conflict: the former is the stationary distribution of the
\emph{full} (unrestricted) mod-$3$ residue chain on the three classes
$\{0, 1, 2\}$, which assigns positive mass to class~$0$ by counting the
density of multiples of $3$ among all integers. The present paper
exclusively uses the active subspace $\mathcal{S}_p$, on which T1
prohibits class~$0$ and the stationary distribution is uniform. All
mutual-information statements below refer to this active subspace.

\subsection{Two measures: Ruelle vs.\ empirical}
\label{ssec:two_measures}

We now introduce the two stationary measures whose discrepancy is the
subject of the paper.

\begin{definition}[Ruelle measure]
\label{def:pi_ruelle}
The \emph{Ruelle measure} of $T_m$ is the unique left probability
eigenvector of $T_m$ for the Perron eigenvalue $\lambda_0 = 1$,
normalised to $\sum_{\mathbf{r}} \pi_m^{\mathrm{Ruelle}}(\mathbf{r}) = 1$:
\begin{equation}
  \pi_m^{\mathrm{Ruelle}}\, T_m \;=\; \pi_m^{\mathrm{Ruelle}},
  \qquad
  \pi_m^{\mathrm{Ruelle}}(\mathbf{r}) \;\geq\; 0
  \;\;\text{for all } \mathbf{r} \in \mathcal{S}_m.
  \label{eq:pi_ruelle_def}
\end{equation}
\end{definition}

\noindent
For $p = 3$, Eq.~\eqref{eq:T3_explicit} gives
$\pi_3^{\mathrm{Ruelle}} = (\tfrac12, \tfrac12)$. For $p \geq 5$,
Eq.~\eqref{eq:Tp_explicit} gives
$\pi_p^{\mathrm{Ruelle}} = \bigl(\tfrac{1}{p-1}, \ldots, \tfrac{1}{p-1}\bigr)$,
uniform on the $(p-1)$ active classes. As explained in
Section~\ref{sec:crt_factor}, the tensor structure of $T_m$ propagates
to $\pi_m^{\mathrm{Ruelle}}$, which factorises as
$\pi_m^{\mathrm{Ruelle}} = \bigotimes_i \pi_{p_i}^{\mathrm{Ruelle}}$.
This factorisation is the algebraic content of Theorem~4.2: under
$\pi_m^{\mathrm{Ruelle}}$, the modular channels are exactly
informationally independent.

The Ruelle measure should be thought of as the \emph{idealised}
stationary measure obtained by treating $T_m$ as a single irreducible
Markov chain and forgetting the arithmetic origin of the trajectory.
In the language of the monograph, it is also the unique
maximum-entropy measure on the set of all admissible trajectories
under the rule $T_{rr} = 0$ for $r \in \mathcal{S}_p$. Both
characterisations are equivalent.

\begin{definition}[Empirical measure]
\label{def:pi_emp}
The \emph{empirical stationary measure} on $\mathbb{Z}/m\mathbb{Z}$ is
defined, for trajectories of the true gap field $\{g_n\}$, by the
long-time limit
\begin{equation}
  \pi_m^{\mathrm{emp}}(\mathbf{r})
  \;:=\;
  \lim_{N \to \infty}
  \frac{1}{N} \sum_{n=1}^{N}
  \mathbbm{1}\!\bigl\{\mathbf{r}_n^{(m)} = \mathbf{r}\bigr\}
  \qquad (\mathbf{r} \in \mathcal{S}_m),
  \label{eq:pi_emp_def}
\end{equation}
whenever the limit exists.
\end{definition}

\noindent
Existence and shift-invariance of the limit~\eqref{eq:pi_emp_def} for
the canonical trajectory $\omega^\ast$ follow from the spectral-gap
estimate of Theorem~T4 in~\cite{PT_Monograph} combined with
Mertens' theorem; we will quote the resulting ergodicity in
Section~\ref{sec:ergodic}, Proposition~\ref{prop:ergodic}.

\begin{remark}[Marginal agreement, joint discrepancy]
\label{rmk:marginal_vs_joint}
The two measures agree at the level of \emph{single-prime marginals}
to numerical precision: for each active prime $p \in \{3, 5, 7\}$,
\begin{equation}
  \bigl\|\, \pi_p^{\mathrm{Ruelle}} - \pi_p^{\mathrm{emp}}\,\bigr\|_{\mathrm{TV}}
  \;<\; 3.5 \times 10^{-7}
  \qquad
  \text{(monograph T3 corpus)}.
  \label{eq:marginal_agree}
\end{equation}
They disagree, however, at the level of \emph{joint} distributions on
$\mathcal{S}_p \times \mathcal{S}_q$ for distinct active $p, q$. The
multi-modular predictor of~\cite[Ch.~math\_structures]{PT_Monograph}
records this disagreement as a substantial mutual information,
$\mathrm{MI}(\mathbf{r}^{(3)}, \mathbf{r}^{(5)}, \mathbf{r}^{(7)})
\sim 1$~nat in the aggregate, and the corpus values
\begin{equation}
  \mathrm{MI}_{\pi^{\mathrm{emp}}}(r^{(3)}\, ; \, r^{(5)}) \;=\; 0.138 \text{ bits},
  \qquad
  \mathrm{MI}_{\pi^{\mathrm{emp}}}(r^{(3)}\, ; \, r^{(7)}) \;=\; 0.283 \text{ bits}
  \label{eq:MI_corpus}
\end{equation}
at the pairwise level (monograph ch.~9, S15.6.259). These two facts ---
marginal agreement~\eqref{eq:marginal_agree} and joint
discrepancy~\eqref{eq:MI_corpus} --- are not contradictory: a joint
distribution that is not a product can have marginals indistinguishable
from those of the product. The corpus values~\eqref{eq:MI_corpus} are
the central data that the closed-form Corollary~\ref{cor:closed_form}
must reproduce.
\end{remark}

\begin{remark}[Where the residual lives]
\label{rmk:residual_arithmetic}
The discrepancy
$\mathcal{I}_{p,q}
:= \mathrm{MI}_{\pi^{\mathrm{emp}}}(r^{(p)}; r^{(q)})
 - \mathrm{MI}_{\pi^{\mathrm{Ruelle}}}(r^{(p)}; r^{(q)})$
is therefore captured entirely by the difference between the two
\emph{joint} laws on $\mathcal{S}_p \times \mathcal{S}_q$, and is
purely arithmetic in origin: it reflects constraints on consecutive
prime gaps that survive once one passes from the idealised Markov chain
$T_m$ to the actual gap field $\{g_n\}$. In particular it has nothing
to do with the geometry~$G(\mu)$ introduced
below.
\end{remark}

\subsection{The emergent geometry $G(\mu)$}
\label{ssec:emergent_geometry}

The Bianchi-I metric of Persistence Theory is reconstructed from the
holonomy data of $T_m$ as follows~\cite[Ch.~10--12]{PT_Monograph}. For
each active prime $p$ and a real parameter $\mu > 0$, set
\begin{align}
  \delta_p(\mu) &\;:=\; \frac{1 - q(\mu)^p}{p},
  &
  \sinp{p}(\mu) &\;:=\; \delta_p(\mu)\bigl(2 - \delta_p(\mu)\bigr),
  \label{eq:holonomy_def}
\\
  \gamp{p}(\mu) &\;:=\;
   -\,\frac{\partial \ln \sinp{p}(\mu)}{\partial \ln \mu},
  &
  a_p(\mu) &\;:=\; \frac{\gamp{p}(\mu)}{\mu},
  \label{eq:scale_factor_def}
\end{align}
where $q(\mu)$ denotes either $\qplus = 1 - 2/\mu$ (couplage branch) or
$q_- = e^{-1/\mu}$ (geometric branch); throughout this paper we take
$q = \qplus$ unless explicitly stated otherwise. The first identity
in~\eqref{eq:holonomy_def} is the holonomy theorem
T6~\cite[Ch.~6]{PT_Monograph}; the anomalous dimension $\gamp{p}$ in
the second line is, up to sign, the logarithmic running of $\sinp{p}$
with respect to the scale~$\mu$.

\begin{definition}[Bianchi-I geometry]
\label{def:G_mu}
The emergent Bianchi-I geometry attached to the active primes
$\{3, 5, 7\}$ is the one-parameter family of diagonal
$(3+1)$-dimensional Lorentzian metrics
\begin{equation}
  G(\mu) \;:=\; \mathrm{diag}\bigl(-1,\; a_3(\mu)^2,\; a_5(\mu)^2,\; a_7(\mu)^2\bigr),
  \qquad \mu \in [\mu_c, \infty),
  \label{eq:G_def}
\end{equation}
where the lower endpoint $\mu_c \approx 6.97$ is the Hartle--Hawking
transition value at which $g_{00} = -\partial^2 \ln \alpha/\partial \mu^2$
changes sign~\cite[Ch.~12]{PT_Monograph}, and the canonical fixed point
of the family is $\mu^\ast = 15$~\cite[Thm.~T5]{PT_Monograph}.
\end{definition}

\noindent
The crucial structural fact, used repeatedly below, is that every
component of $G(\mu)$ depends only on the spectrum of $T_m$ through the
holonomy identities~\eqref{eq:holonomy_def}--\eqref{eq:scale_factor_def}:
the Bianchi-I geometry is, in a precise measure-theoretic sense, a
\emph{function of the spectrum} of the transfer matrix, not of any
particular trajectory of the gap field. We make this precise in
Section~\ref{sec:geometric_fiber} (Lemma~\ref{lem:G_measurable}); it is
the key technical ingredient of Theorem~\ref{thm:empirical_invariance}.

\subsection{Summary table of objects}
\label{ssec:summary_objects}

Table~\ref{tab:setup_glossary} collects the objects introduced above
in the order in which they will be used in the proofs. We adopt the
notation throughout the paper without further comment.

\begin{table}[h]
\centering
\caption{Objects defined in Section~\ref{sec:setup} and their nature.
``Spectral'' means: a measurable functional of $\mathrm{Spec}(T_m)$
alone; ``trajectory'' means: a functional of $\omega^\ast \in \Omega$;
``arithmetic'' means: fixed by the sieve algorithm, independent
of~$\mu$.}
\label{tab:setup_glossary}
\renewcommand{\arraystretch}{1.15}
\begin{tabular}{@{}l l l@{}}
\toprule
Symbol & Definition & Nature \\
\midrule
$g_n$
 & $p_{n+1} - p_n$ (Eq.~\eqref{eq:gap_def})
 & arithmetic \\
$r_n^{(m)}$
 & $g_n \bmod m$ (Eq.~\eqref{eq:residue_def})
 & arithmetic \\
$\mathcal{A}_p$
 & $\sigma$-algebra generated by $\{r_n^{(p)}\}_n$ (Eq.~\eqref{eq:Ap_def})
 & arithmetic \\
$\mathcal{S}_p$
 & active subspace $\{1, \ldots, p-1\}$
 & arithmetic \\
$T_p$
 & restricted transfer matrix on $\mathcal{S}_p$
   (Eqs.~\eqref{eq:T3_explicit},~\eqref{eq:Tp_explicit})
 & arithmetic \\
$T_m$
 & tensor product $\bigotimes_{p \mid m} T_p$
   (Eq.~\eqref{eq:Tm_tensor_preview})
 & arithmetic \\
$\pi_m^{\mathrm{Ruelle}}$
 & left Perron eigenvector of $T_m$ (Def.~\ref{def:pi_ruelle})
 & spectral \\
$\pi_m^{\mathrm{emp}}$
 & empirical occupation of $\{r_n^{(m)}\}_n$ (Def.~\ref{def:pi_emp})
 & trajectory \\
$\sinp{p}(\mu)$
 & $\delta_p(2 - \delta_p)$, $\delta_p = (1 - q^p)/p$
   (Eq.~\eqref{eq:holonomy_def})
 & spectral \\
$\gamp{p}(\mu)$
 & $-\partial \ln \sinp{p}/\partial \ln \mu$
   (Eq.~\eqref{eq:scale_factor_def})
 & spectral \\
$a_p(\mu)$
 & $\gamp{p}(\mu)/\mu$ (Eq.~\eqref{eq:scale_factor_def})
 & spectral \\
$G(\mu)$
 & $\mathrm{diag}(-1, a_3^2, a_5^2, a_7^2)$ (Eq.~\eqref{eq:G_def})
 & spectral \\
$\Phi_p$
 & bounded $\mathcal{A}_p$-measurable observable
 & trajectory \\
\bottomrule
\end{tabular}
\end{table}

The classification ``spectral~vs.~trajectory'' in the right column will
do most of the work in the proof of
Theorem~\ref{thm:empirical_invariance}: spectral objects are
$\mathcal{T}$-invariant under the shift and hence
$\pi_m^{\mathrm{emp}}$-almost-surely constant by ergodicity, while
trajectory observables carry the mutual information. The proof reduces
to checking that $G(\mu)$ is on the spectral side of this divide for
every $\mu \in [\mu_c, \infty)$.
